\subsection{PARACOMPACTNESS OF METRIZABLE SPACES}

The following theorem sharpens the result of Proposition 2 of no. 1:

THEOREM 4. Every metrizable space is paracompact.

The theorem is a consequence of the following four lemmas.

LEMMA 3. Let \(\Re = (U_{\alpha})_{\alpha \in A}\) be an open covering of a metrizable space \(X\) . Then there is a sequence \((\mathfrak{S}_n)\) of locally finite families of open subsets of \(X\) , such that \(\mathfrak{S} = \bigcup \mathfrak{S}_n\) is an open covering of \(X\) which refines \(\Re\) .

Let \(d\) be a metric on \(\mathbf{X}\) compatible with its topology. For each \(\alpha \in \mathbf{A}\) and each integer \(n\) , let \(\mathbf{F}_{n\alpha}\) denote the set of all \(x \in \mathbf{U}_{\alpha}\) such that \(d(x, \mathbf{X} - \mathbf{U}_{\alpha}) \geqslant 2^{-n}\) . Since \(\mathbf{X} - \mathbf{U}_{\alpha}\) is closed, we have \(\mathbf{U}_{\alpha} = \bigcup_{n} \mathbf{F}_{n\alpha}\) . Well-order the set \(\mathbf{A}\) ; for each \(\alpha \in \mathbf{A}\) and each integer \(n\) , let \(\mathbf{G}_{n\alpha}\) be the set of all \(x \in \mathbf{F}_{n\alpha}\) such that \(x \notin \mathbf{F}_{n+1,\beta}\) for all \(\beta < \alpha\) , and let \(\mathbf{V}_{n\alpha}\) be the set of all \(y \in \mathbf{X}\) such that \(d(y, \mathbf{G}_{n\alpha}) - 2^{-n-3}\) . \(\mathbf{V}_{n\alpha}\) is clearly an open set; on the other hand, \(\mathbf{V}_{n\alpha} \subset \mathbf{U}_{\alpha}\) , because for each \(y \in \mathbf{V}_{n\alpha}\) there exists \(x \in \mathbf{G}_{n\alpha}\) such that \(d(x, y) \leqslant 2^{-n-1}\) , and since \(x \in \mathbf{F}_{n\alpha}\) , we have

\[
d (y, \mathrm{X} - \mathrm{U} _ {\alpha}) \geqslant d (x, \mathrm{X} - \mathrm{U} _ {\alpha}) - d (x, y) \geqslant 2 ^ {- n - 1},
\]

so that \(y \in \mathbf{U}_{\alpha}\) . Furthermore, for each \(x \in \mathbf{X}\) let \(\alpha\) be the smallest index in \(\mathbf{A}\) such that \(x \in \mathbf{U}_{\alpha}\) ; then there is an integer \(n\) such that \(x \in \mathrm{F}_{n\alpha}\) , and the definition of \(\alpha\) implies that \(x \in \mathrm{G}_{n\alpha}\) , so that \(x \in \mathrm{V}_{n\alpha}\) . This shows that if we put \(\mathfrak{S}_n = (\mathrm{V}_{n\alpha})_{\alpha \in \mathbf{A}}\) , then \(\mathfrak{S} = \bigcup_{n} \mathfrak{S}_n\) is an open

covering of X which refines \(\Re\) ; thus it remains to be shown that each of the families \(\mathfrak{S}_n\) is locally finite. To this end we shall first show that \(d(\mathrm{G}_{n\alpha}, \mathrm{G}_{n\beta}) \geqslant 2^{-n-1}\) if \(\alpha \neq \beta\) . Suppose that \(\beta < \alpha\) ; then if \(x \in \mathrm{G}_{n\alpha}\) and \(y \in \mathrm{F}_{n\beta}\) we have \(x \notin \mathrm{F}_{n+1, \beta}\) by definition, hence \(d(x, \mathrm{X} - \mathrm{U}_{\beta}) < 2^{-n-1}\) and \(d(y, \mathrm{X} - \mathrm{U}_{\beta}) \geqslant 2^{-n}\) , and therefore \(d(x, y) \geqslant 2^{-n-1}\) ; since \(\mathrm{G}_{n\beta} \subset \mathrm{F}_{n\beta}\) , the assertion follows. From this it follows immediately, using the definition of \(\mathrm{V}_{n\alpha}\) and \(\mathrm{V}_{n\beta}\) , that \(d(\mathrm{V}_{n\alpha}, \mathrm{V}_{n\beta}) \geqslant 2^{-n-2}\) . From this last inequality we deduce that, for each \(z \in \mathrm{X}\) , the open ball with centre \(z\) and radius \(2^{-n-3}\) meets at most one set of the family \(\mathfrak{S}_n\) ; hence \(\mathfrak{S}_n\) is a locally finite family, and the proof is complete.

LEMMA 4. Let \((\mathfrak{S}_n)\) be a sequence of locally finite families of open sets in a topological space \(\mathbf{X}\) , such that

\[
\mathfrak {S} = \bigcup_ {n} \mathfrak {S} _ {n}
\]

is a covering of X. Then there exists a locally finite (but not necessarily open) covering B of X which refines S.

Let \(\mathbf{E}_n\) be the open set in \(\mathbf{X}\) which is the union of all the sets of \(\mathfrak{S}_n\) ;

let \(U_{n}\) denote \(\bigcup_{k=1}^{n}E_{k}\) and let \(A_{n}\) denote \(U_{n}-U_{n-1}\) (with \(U_{0}=\varnothing\) ). Consider the set B of subsets \(V\cap A_{n}\) , where \(V\in S_{n}\) and n is any integer; we shall show that B satisfies the conditions of the lemma. For each \(x\in X\) there is an integer n such that \(x\in A_{n}\) , since the \(A_{n}\) form a partition of X; thus \(x\in E_{n}\) and there exists \(V\in S_{n}\) such that \(x\in V\) ; so \(x\in V\cap A_{n}\) , and we have proved that B is a covering of X. Clearly, this covering refines S. On the other hand, for each \(x\in X\) there exists an integer n such that \(x\in U_{n}\) ; since \(U_{n}\) is open and the \(S_{m}\) are locally finite families, there exists a neighbourhood \(W_{m}\) of x, for each m, contained in \(U_{n}\) , which meets only a finite number of sets of \(S_{m}\) ; hence the neighbourhood \(W=\bigcap_{m=1}^{n}W_{m}\) of x meets only a finite number of sets of B, because \(W\cap A_{p}=\varnothing\) for p\textgreater n.~Hence B is locally finite, and Lemma 4 is proved.

LEMMA 5. Let X be a regular space such that, for each open covering R of X, there exists a (not necessarily open) locally finite covering A of X which refines R. Then for each open covering R of X there exists a locally finite closed covering F of X which refines R.

Let \(\mathfrak{R}\) be any open covering of X. For each \(x \in \mathbf{X}\) there is an open set \(\mathbf{U} \in \mathfrak{R}\) which contains \(x\) , and therefore (since X is regular) an open neighbourhood \(\mathbf{V}_x\) of \(x\) such that \(\overline{\mathbf{V}}_x \subset \mathbf{U}\) . The family \(\mathfrak{B}\) formed by the \(\mathbf{V}_x\) is an open covering of X, hence by hypothesis there is a locally finite covering \(\mathfrak{B}\) which is finer than \(\mathfrak{B}\) . Let \(\mathfrak{F}\) be the family of closures of the sets of \(\mathfrak{B}\) . Since the covering \(\mathfrak{B}'\) formed by the \(\overline{\mathbf{V}}_x\) is finer than \(\mathfrak{R}\) , and since \(\mathfrak{F}\) is finer than \(\mathfrak{B}'\) , it follows that \(\mathfrak{F}\) is a closed covering of X which refines \(\mathfrak{R}\) . But also \(\mathfrak{F}\) is locally finite, because if an open set does not meet a set \(\mathbf{B} \in \mathfrak{B}\) , then it does not meet its closure \(\overline{\mathbf{B}}\) either.

LEMMA 6. Let X be a Hausdorff space such that, given any open covering R of X, there exists a locally finite closed covering F of X which refines R. Then X is paracompact.

Let \(\mathfrak{R}\) be any open covering of X. We have to show that there is a locally finite open covering of X which refines \(\mathfrak{R}\) . Let \(\mathfrak{A}\) be a locally finite covering (closed or not) of X which refines \(\mathfrak{R}\) ; for each \(x \in \mathbf{X}\) , let \(W_x\) be an open neighbourhood of \(x\) which meets only a finite number of sets of \(\mathfrak{A}\) . The family \(\mathfrak{B}\) of sets \(W_x\) is an open covering of X. Let \(\mathfrak{F}\) be a locally finite closed covering of X which refines \(\mathfrak{B}\) . For each \(A \in \mathfrak{A}\) , let \(U_A\) be a set of \(\mathfrak{R}\) which contains A, and let \(C_A\) be the union of the sets \(F \in \mathfrak{F}\) such that \(A \cap F = \emptyset\) . Since \(\mathfrak{F}\) is locally finite, \(C_A\) is closed in X (Chapter I, § 1, no. 6, Proposition 4) and therefore \(A' = U_A \cap (X - C_A)\) is open. Since we have \(A \cap C_A = \emptyset\) and \(A \subset U_A\) , it follows that \(A \subset A'\) , and the family \(\mathfrak{A}'\) of sets \(A'\) , as A runs through \(\mathfrak{A}\) , is an open covering of X; moreover, since \(A' \subset U_A \in \mathfrak{R}\) , \(\mathfrak{A}'\) refines \(\mathfrak{R}\) . It remains to show that \(\mathfrak{A}'\) is locally finite. For each \(x \in X\) there is a neighbourhood T of \(x\) which meets only a finite number of sets of \(\mathfrak{F}\) , say \(F_1, \ldots, F_n\) . Since each \(F_i\) is contained in a set of the form \(W_{y_i}\) , by definition \(F_i\) meets only a finite number of sets of \(\mathfrak{A}\) ; let \(A_{ij} (i \leqslant j \leqslant s_i)\) be these sets. If A is a set of \(\mathfrak{A}\) other than one of the \(A_{ij} (i \leqslant i \leqslant n, i \leqslant j \leqslant s_i)\) , it follows from the definitions that \(A'\) meets none of the \(F_i\) , and therefore does not meet \(T \subset \bigcup_{i=1}^{n} F_i\) . This completes the proof of Lemma 6 and of Theorem 4.

